\documentclass[UTF8,11pt]{ctexart}
\usepackage[a4paper,margin=24mm]{geometry}
\usepackage{amsmath,amssymb,amsthm,mathtools,mathrsfs}
\usepackage[all]{xy}
\usepackage{xurl,hyperref,enumitem}
\usepackage{tikz}
\usetikzlibrary{arrows.meta,cd,decorations.markings,calc,patterns}
\hypersetup{hidelinks,breaklinks=true}
\setlength{\emergencystretch}{3em}
\allowdisplaybreaks
\newtheorem*{theorem}{Theorem}
\newtheorem*{thm}{Theorem}
\newtheorem*{lemma}{Lemma}
\newtheorem*{proposition}{Proposition}
\newtheorem*{prop}{Proposition}
\newtheorem*{corollary}{Corollary}
\newtheorem*{cor}{Corollary}
\newtheorem*{claim}{Claim}
\theoremstyle{definition}
\newtheorem*{definition}{Definition}
\newtheorem*{defn}{Definition}
\newtheorem*{remark}{Remark}
\newtheorem*{example}{Example}
\newcommand{\R}{\mathbb R}
\newcommand{\C}{\mathbb C}
\newcommand{\Z}{\mathbb Z}
\newcommand{\Q}{\mathbb Q}
\newcommand{\N}{\mathbb N}
\newcommand{\F}{\mathbb F}
\newcommand{\D}{\mathbb D}
\newcommand{\bB}{\mathbb B}
\newcommand{\bC}{\mathbb C}
\newcommand{\bR}{\mathbb R}
\newcommand{\bZ}{\mathbb Z}
\newcommand{\bN}{\mathbb N}
\newcommand{\bQ}{\mathbb Q}
\newcommand{\bD}{\mathbb D}
\newcommand{\bF}{\mathbb F}
\newcommand{\CP}{\mathbb{CP}}
\newcommand{\RP}{\mathbb{RP}}
\newcommand{\sO}{\mathscr O}
\newcommand{\sI}{\mathscr I}
\newcommand{\sF}{\mathscr F}
\newcommand{\sG}{\mathscr G}
\newcommand{\sH}{\mathscr H}
\newcommand{\sR}{\mathscr R}
\newcommand{\sM}{\mathscr M}
\newcommand{\sV}{\mathscr V}
\newcommand{\sU}{\mathscr U}
\DeclareMathOperator{\Hom}{Hom}
\DeclareMathOperator{\Ext}{Ext}
\DeclareMathOperator{\Tor}{Tor}
\DeclareMathOperator{\Spec}{Spec}
\DeclareMathOperator{\Proj}{Proj}
\DeclareMathOperator{\supp}{supp}
\DeclareMathOperator{\im}{im}
\DeclareMathOperator{\id}{id}
\DeclareMathOperator{\Tr}{Tr}
\DeclareMathOperator{\tr}{tr}
\DeclareMathOperator{\rank}{rank}
\DeclareMathOperator{\ord}{ord}
\DeclareMathOperator{\dist}{dist}
\DeclareMathOperator{\End}{End}
\DeclareMathOperator{\Aut}{Aut}
\DeclareMathOperator{\Gal}{Gal}
\DeclareMathOperator{\vol}{vol}
\DeclareMathOperator{\Cl}{Cl}
\DeclareMathOperator{\coker}{coker}
\DeclareMathOperator{\codim}{codim}
\DeclareMathOperator{\divergence}{div}
\DeclareMathOperator{\Ric}{Ric}
\DeclareMathOperator{\grad}{grad}
\DeclareMathOperator{\Hess}{Hess}
\newcommand{\abs}[1]{\left\lvert#1\right\rvert}
\newcommand{\sabs}[1]{\lvert#1\rvert}
\newcommand{\babs}[1]{\bigl\lvert#1\bigr\rvert}
\newcommand{\Babs}[1]{\Bigl\lvert#1\Bigr\rvert}
\newcommand{\bbabs}[1]{\biggl\lvert#1\biggr\rvert}
\newcommand{\BBabs}[1]{\Biggl\lvert#1\Biggr\rvert}
\newcommand{\norm}[1]{\left\lVert#1\right\rVert}
\newcommand{\snorm}[1]{\lVert#1\rVert}
\newcommand{\bnorm}[1]{\bigl\lVert#1\bigr\rVert}
\newcommand{\Bnorm}[1]{\Bigl\lVert#1\Bigr\rVert}
\newcommand{\bbnorm}[1]{\biggl\lVert#1\biggr\rVert}
\newcommand{\BBnorm}[1]{\Biggl\lVert#1\Biggr\rVert}
\newcommand{\widebar}[1]{\overline{#1}}
\newcommand{\nicefrac}[2]{\frac{#1}{#2}}
\newcommand{\linnprod}[2]{\langle#1,#2\rangle}
\newcommand{\myindex}[1]{#1}
\newcommand{\glsadd}[1]{}
\newcommand{\avoidbreak}{}
\newcommand{\sourceref}[1]{\textnormal{[原文：\nolinkurl{#1}]}}
\newcommand{\thmref}[1]{\sourceref{#1}}
\newcommand{\lemmaref}[1]{\sourceref{#1}}
\newcommand{\propref}[1]{\sourceref{#1}}
\newcommand{\corref}[1]{\sourceref{#1}}
\newcommand{\defnref}[1]{\sourceref{#1}}
\newcommand{\exampleref}[1]{\sourceref{#1}}
\newcommand{\exerciseref}[1]{\sourceref{#1}}
\newcommand{\figureref}[1]{\sourceref{#1}}
\newcommand{\sectionref}[1]{\sourceref{#1}}
\newcommand{\appendixref}[1]{\sourceref{#1}}
\newcommand{\stacksref}[2]{\href{https://stacks.math.columbia.edu/tag/#1}{#2 [#1]}}


\newcommand{\E}{\mathbb E}
\newcommand{\Prob}{\mathbb P}
\newcommand{\Reff}{\mathcal R_{\mathrm{eff}}}
\newcommand{\eps}{\varepsilon}
\DeclareMathOperator{\diam}{diam}
\DeclareMathOperator{\Var}{Var}
\DeclareMathOperator{\Crit}{Crit}
\newcommand{\dd}{\,\mathrm d}


\begin{document}
\section*{090\quad Lie群闭子群的Cartan光滑结构定理}
\subsection*{来源与整理范围}
Terence Tao，254A, Notes 2: Building Lie structure from representations and metrics，2011-09-08；Theorem 2、Section 1 的线性子空间论证、Lemmas 10--11 与主定理结论。
\par\url{https://terrytao.wordpress.com/2011/09/08/254a-notes-2-building-lie-structure-from-representations-and-metrics/}
\par\textbf{恢复方式（整理者说明）：}依据人类原文重新整理以补齐缺失题号；并非旧题证文件的逐字节恢复；2026-09-28。
\subsection*{作者命题与人类证明}

\begin{theorem}[Cartan's closed-subgroup theorem; Tao's Theorem 2]
A closed subgroup $H$ of a finite-dimensional real Lie group $G$ is a
smooth embedded submanifold of $G$, and hence is a Lie group with its
subspace topology.
\end{theorem}
\paragraph{One-parameter subgroups.}
A one-parameter subgroup is a continuous homomorphism $\R\to G$.
For a Lie group these are exactly the maps $t\mapsto\exp(tX)$,
$X\in\mathfrak g$. Thus the one-parameter subgroups of $H$ identify with
\[
 \mathfrak h=\{X\in\mathfrak g:\exp(tX)\in H\text{ for all }t\in\R\}.
\]
This set contains zero and is closed under real scalar multiplication.
If $X,Y\in\mathfrak h$, the Lie product formula gives
\[
 \exp(t(X+Y))=
 \lim_{n\to\infty}\bigl(\exp(tX/2^n)\exp(tY/2^n)\bigr)^{2^n}.
\]
Every term lies in $H$, so its limit lies in $H$ by closedness.
Consequently $\mathfrak h$ is a linear subspace of $\mathfrak g$.

\begin{lemma}[Nontrivial tangent direction; Tao's Lemma 10]
If the identity is not isolated in $H$, then $\mathfrak h\ne\{0\}$.
\end{lemma}
\begin{proof}
Choose $h_n\in H\setminus\{e\}$ with $h_n\to e$.
The exponential map is a local diffeomorphism near zero, so for large $n$
write $h_n=\exp(X_n)$ with $X_n\ne0$, $X_n\to0$.
Fix a norm on the finite-dimensional vector space $\mathfrak g$.
By compactness of its unit sphere, after passing to a subsequence
\[
 X_n/\|X_n\|\longrightarrow\omega,\qquad\|\omega\|=1.
\]
For $t>0$,
$\lfloor t/\|X_n\|\rfloor X_n\to t\omega$ and therefore
\[
 h_n^{\lfloor t/\|X_n\|\rfloor}
 =\exp\bigl(\lfloor t/\|X_n\|\rfloor X_n\bigr)
 \longrightarrow\exp(t\omega).
\]
The powers belong to $H$; their limits do too. Taking inverses also covers
$t<0$, so $\omega\in\mathfrak h$ and $\mathfrak h$ is nontrivial.
\end{proof}

\begin{lemma}[Local exponential chart; Tao's Lemma 11]
There are identity and zero neighbourhoods $U\subset H$ and
$V\subset\mathfrak h$ for which $\exp:V\to U$ is a homeomorphism.
\end{lemma}
\begin{proof}
Take a small neighbourhood $V$ of zero in $\mathfrak h$ on which the
exponential map is injective and is a homeomorphism onto its image.
If $\exp(V)$ contains a neighbourhood of the identity in $H$, the claim
follows after restricting neighbourhoods. Otherwise there are
$h_n\in H\setminus\exp(V)$ with $h_n\to e$. Write
$h_n=\exp(X_n)$, where $X_n\to0$ in $\mathfrak g$.

Choose a vector-space complement $\mathfrak k$ of $\mathfrak h$:
\[
 \mathfrak g=\mathfrak h\oplus\mathfrak k.
\]
No Lie-algebra condition is imposed on $\mathfrak k$. By the inverse
function theorem, the smooth map
\[
 \mathfrak h\times\mathfrak k\longrightarrow G,
 \qquad(Y,Z)\longmapsto\exp(Y)\exp(Z)
\]
is a local diffeomorphism at $(0,0)$. For all large $n$ we may write
\[
 h_n=\exp(Y_n)\exp(Z_n),\qquad
 Y_n\in\mathfrak h,\quad Z_n\in\mathfrak k,\quad Y_n,Z_n\to0.
\]
Since $h_n\notin\exp(V)$ and $Y_n$ is eventually in $V$, we have $Z_n\ne0$.
Choose a norm on $\mathfrak k$ and pass to a subsequence such that
$Z_n/\|Z_n\|\to\omega$ with $\|\omega\|=1$. In particular,
$\omega\in\mathfrak k\setminus\mathfrak h$.

But $\exp(Z_n)=\exp(-Y_n)h_n\in H$. The integer-power argument in the
previous lemma now implies $\exp(t\omega)\in H$ for every $t\in\R$.
Thus $\omega\in\mathfrak h$, a contradiction.
\end{proof}

\begin{proof}[Conclusion of Cartan's theorem]
The preceding lemma shows that $H$ agrees locally near the identity with
$\exp(V)$, where $V$ is open in the linear subspace $\mathfrak h$.
The ambient exponential chart is smooth with smooth inverse, so this is
a smooth embedded submanifold near the identity. Translation by elements
of $H$ gives such charts at every point of $H$. The group operations are
restrictions of the smooth operations of $G$, hence are smooth in this
submanifold structure. Therefore $H$ is a Lie group.
\end{proof}

\subsection*{整理说明与先修范围（非作者正文）}
按作者原证明转录，保留切向非零极限、横截补空间和排除额外小元素的核心步骤。先修为环境 Lie 群中的指数映射局部微分同胚、单参数子群与 Lie 代数的对应、Lie 乘积公式和有限维逆函数定理。这里不预先假定 $H$ 有微分结构，也不引用闭子群定理自身。

本题是全局闭子群定理，不把作者另留作练习的局部 Cartan 或 von Neumann 定理并入题面。网页的局部 homeomorphism 在使用光滑逆函数定理的地方写清为 diffeomorphism，未更换证明。难点在于把闭性的序列极限转换为实际单参数子群，并证明这些方向已穷尽闭子群的局部结构。
\subsection*{可修改的简短标注（非作者正文）}
$c$：拓扑闭子群的光滑流形结构

$a$：单参数子群；指数映射；Lie乘积公式；紧单位球面；横截补空间；逆函数定理

\texttt{cross\_theory\_status = yes}。将闭子群中的小元素及其幂缩放转成Lie代数方向，再用横截分解排除额外方向并返回子流形坐标；见两条引理。

全部标注为非穷尽、可修改初稿。
\subsection*{来源许可与编辑归属}
作者 About 页 Copyright etc. 允许带来源引用复用合理篇幅（例如单篇文章）；本题为原文指定部分的编辑转录。许可入口：https://terrytao.wordpress.com/about/ 。
ChatGPT 加入题号、中文来源说明、整理范围和初稿标注；编辑说明与人类证明分列。
\end{document}
